\documentclass[11pt]{scrartcl}
\usepackage[a4paper,margin=2.4cm]{geometry}
\usepackage{../common/rlstudy}

\title{\textbf{Week 23 — Extended reading}\\\large Shielding, CBFs, Domain Randomisation, Robust RL}
\author{STUDY-RL}
\date{\today}

\begin{document}
\maketitle

\section*{Paper 1 --- ``Safe Reinforcement Learning via Shielding''}
\textbf{Authors:} Mohammed Alshiekh, Roderick Bloem, Rüdiger Ehlers, Bettina
Könighofer, Scott Niekum, Ufuk Topcu. \textbf{Year:} 2018. \textbf{Venue:} AAAI.
\textbf{arXiv:} 1708.08611.

\paragraph{Summary.} Introduces the concept of an RL \emph{shield}: a reactive
correct-by-construction component synthesised from an $\omega$-regular safety
specification, that overrides the agent's action whenever it would lead to a
violation. Both pre-shield (filters action set) and post-shield (corrects
chosen action) variants are described. Soundness is proved against the
specification.

\paragraph{Why it matters.} The canonical reference for the formal-methods
side of safe RL. Our \texttt{project/shield.py} is a clean implementation of
the pre-shield variant for the running $\mathbf{G}\,\neg$danger specification.

\paragraph{Annotated key passage.}
\begin{quote}
``A correct-by-construction shield is a reactive system that monitors the
actions chosen by the learning agent and corrects them if necessary, such
that the temporal logic specification is satisfied. The shield is computed
offline from the specification and a (possibly abstract) model of the
environment.''
\end{quote}
\emph{Commentary.} The key conceptual move is the \emph{decoupling} of safety
and learning. The shield's soundness does not depend on the learner; the
learner's correctness does not depend on the shield. This factoring is what
makes the approach scale to deep RL: you can plug a shield in front of any
policy.

% -------------------------------------------------------------------
\section*{Paper 2 --- ``Safe Exploration in Continuous Action Spaces''}
\textbf{Authors:} Gal Dalal, Krishnamurthy Dvijotham, Matej Vecerik, Todd Hester,
Cosmin Paduraru, Yuval Tassa. \textbf{Year:} 2018. \textbf{arXiv:} 1801.08757.

\paragraph{Summary.} Augments an actor with a final \emph{safety layer} ---
a closed-form (or differentiable QP) projection that maps the actor's proposed
action onto the half-space of safe actions defined by a linearised constraint.
The whole pipeline is differentiable so gradients flow through the projection.

\paragraph{Why it matters.} Where shielding is the discrete-action poster
child, the safety layer is the continuous-action poster child. Closed-form
projections for a single linear constraint, OptNet-style QPs for multiple.

\paragraph{Annotated key passage.}
\begin{quote}
``The Safety Layer is a closed-form linear projection $a^\star = a - [g(s) +
\partial g/\partial a \cdot a]_+ / \|\partial g/\partial a\|^2 \cdot \partial
g/\partial a$. Crucially, this projection is differentiable, so the safety
layer can be inserted at training time without breaking the policy
gradient.''
\end{quote}
\emph{Commentary.} The differentiability is the whole point. A non-differentiable
safety wrapper would force you into surrogate gradient estimators; OptNet (Amos
\& Kolter 2017) made the QP itself a primitive that PyTorch can backprop
through.

% -------------------------------------------------------------------
\section*{Paper 3 --- ``Domain Randomization for Transferring Deep Neural Networks''}
\textbf{Authors:} Josh Tobin, Rachel Fong, Alex Ray, Jonas Schneider, Wojciech
Zaremba, Pieter Abbeel. \textbf{Year:} 2017. \textbf{arXiv:} 1703.06907.

\paragraph{Summary.} Trains policies on a wide \emph{distribution} of
simulator parameters (lighting, friction, textures) and transfers zero-shot
to a real robot. The policy must implicitly learn invariances across the
distribution.

\paragraph{Why it matters.} The simplest robust-RL technique. Often the first
thing tried in sim-to-real; frequently \emph{the best} thing tried, because
it dovetails with any underlying RL algorithm.

\paragraph{Annotated key passage.}
\begin{quote}
``The key insight is that, by randomising aspects of the simulator that may
otherwise overfit, the network is forced to learn robust features --- and the
real world simply becomes another sample from the training distribution.''
\end{quote}
\emph{Commentary.} The same principle is data augmentation, lifted to RL.
The critical knob is the \emph{support} of the randomisation: too narrow and
the policy doesn't generalise; too broad and training collapses (no signal).
ADR (OpenAI Rubik's Cube, arXiv:1910.07113) automates this expansion.

% -------------------------------------------------------------------
\section*{Paper 4 --- ``Sim-to-Real Transfer with Foundation Models''}
\textbf{Authors:} Longchao Da, Justin Turnau, Thirulogasankar Pranav Kutralingam,
Alvaro Velasquez, Paulo Shakarian, Hua Wei. \textbf{Year:} 2025.
\textbf{arXiv:} 2502.13187.

\paragraph{Summary.} A survey of how vision--language foundation models are
being plugged into the sim-to-real pipeline: as world models, as observation
encoders, as reward labellers, and as scene generators for DR.

\paragraph{Why it matters.} The 2025 frontier for robust RL. The foundation
model gives a near-universal feature space; DR happens \emph{in feature
space}, not in pixel space, which is much more sample-efficient.

\paragraph{Annotated key passage.}
\begin{quote}
``Foundation-model encoders project simulated and real observations into a
shared semantic space, transforming the sim-to-real gap from a pixel-statistics
problem into a feature-alignment problem. This is the modern descendant of
domain randomisation.''
\end{quote}
\emph{Commentary.} The clean way to view this: pre-2020 DR was about making
the policy invariant; post-2024 DR is about making the encoder invariant.
The policy can be much simpler when the encoder already abstracts away
nuisance variation.

% -------------------------------------------------------------------
\section*{Paper 5 --- ``Learning-Based Shielding under Unknown Dynamics''}
\textbf{Authors:} (Multi-author, see arXiv). \textbf{Year:} 2024.
\textbf{arXiv:} 2410.07359.

\paragraph{Summary.} Extends the shielding framework to settings where the
environment model is learned rather than given. Provides probabilistic
soundness bounds in terms of model uncertainty (PAC-style).

\paragraph{Why it matters.} Lifts the most restrictive assumption of
Alshiekh-style shielding (a known one-step model). Critical for real
deployments where the world is partially observed and partially learned.

\paragraph{Annotated key passage.}
\begin{quote}
``We show that, when the learned dynamics model is calibrated, the shield's
soundness guarantee degrades gracefully: the probability of an unsafe
trajectory is bounded by the model's coverage failure probability plus the
spec's reachability margin.''
\end{quote}
\emph{Commentary.} This is the bridge from formal methods to statistical
learning: shields are no longer ``correct or broken'' but ``correct with
probability $1 - \delta$''. The same language is showing up in conformal
prediction wrappers around RL agents (a 2024 healthcare-RL theme).

\paragraph{Sourced from.} \texttt{research/safety\_rl\_sota.md}, sections 2
(shielding, CBFs) and 4 (robust RL, DR).

\end{document}
